\section{Reproducibility}
\label{app:repro}

\subsection{Software}

The project is written in Python 3.14 and managed with \texttt{uv}, which pins
the environment from a lock file. The simulation and learning stack uses
\texttt{numpy}, \texttt{polars} and \texttt{pyarrow} for data handling,
\texttt{gymnasium} for the environment interface\cite{towers2024gymnasium},
and \texttt{stable-baselines3} for the PPO trainer\cite{raffin2021stablebaselines3}.
The report is built from LaTeX with a wrapper that selects among
\texttt{tectonic}, \texttt{latexmk} and \texttt{pdflatex}.

\subsection{Program inventory}

Table~\ref{tab:app:programs} lists every program and what it computes. The
analysis programs are deterministic given the artefacts they read; the
learning runs are seeded. The repository exposes three console commands (data,
simulator and dashboard), and the reproduction wrapper
\texttt{scripts/report/reproduce.sh} is the single main program: it invokes
the console commands and the analysis programs in order and rebuilds the
report.

\begin{table}[tbp]
\caption{Program inventory. Package programs live under \texttt{src/undertow/},
the analysis programs in \texttt{scripts/report/}, the configurations in
\texttt{configs/}, and the report tooling in \texttt{report/} and
\texttt{scripts/}. The analysis programs write the CSV artefacts that the
appendix tables and the main-body numbers read.}
\label{tab:app:programs}
\centering
\small
\begin{tabular}{@{}L{0.34\linewidth}L{0.60\linewidth}@{}}
\toprule
\textbf{Program} & \textbf{What it computes} \\
\midrule
data package & data pipeline: snapshot construction, quality checks and validation reports \\
simulator package & simulator, exact backtester, parity harness, training and evaluation \\
dashboard package & live inspection of runs, telemetry and rollouts; not used for any reported result \\
configuration files & every experiment configuration (pools, splits, grids, costs, training) \\
\texttt{reproduce.sh} & main reproduction program: runs the stages below in order \\
\texttt{analyze\_main\_run.py} & main joint run: per-window metrics, paired bootstrap intervals, action mix, per-regime tables, decomposition \\
\texttt{static\_widths.py} & static one-shot widths on the training and evaluation splits \\
\texttt{extra\_benchmarks.py} & the two scripted recentring rules on the evaluation windows \\
\texttt{boundary\_sweep.py} & per-window volatility and volume, gas and capital reruns, analytical LVR \\
\texttt{contrast\_fee\_tier.py} & the full ladder on the $0.30\,\%$ pool \\
\texttt{analyze\_regime\_ablation.py} & the regime-feature ablation arm \\
\texttt{make\_appendix\_tables.py} & every table in Appendices~\ref{app:results} and \ref{app:ablations}, from the CSV artefacts \\
\texttt{make\_figures.py} & the four story figures of the main body \\
\texttt{build\_report.sh} & builds the report PDF from the LaTeX sources \\
\texttt{plagiarism\_check.py} & originality screen against the cited papers and the repository's own prose \\
\bottomrule
\end{tabular}
\end{table}

\subsection{Data files}

The two dataset snapshots and the recorded evaluation artefacts are
distributed with the submission archive; they are excluded from version
control because of their size. The primary pool's snapshot is approximately
$812$\,MB and the contrast pool's approximately $293$\,MB, partitioned by
stream and block bucket. Each carries a manifest with row counts and block
ranges and a validation report; Appendix~\ref{app:data} summarises both. The
reproduction program does not re-fetch the data because the fetch step needs
API credentials; every reported number is produced from the distributed
snapshots and the recorded evaluation artefacts under \texttt{results/}.

\subsection{How to run}

From the repository root, the wrapper reproduces the analyses, tables,
figures, PDF and checks:

{\footnotesize
\begin{verbatim}
uv sync
scripts/report/reproduce.sh            # analyses, tables, figures, PDF, checks
scripts/report/reproduce.sh --train    # the above, after retraining both arms
\end{verbatim}
}

The \texttt{--train} stage retrains the learned policy and the ablation arm;
it takes several hours on the hardware below. The analyses take under an hour
in total, and the report build takes seconds.

\subsection{Hardware and runtimes}

The experiments ran on the author's training machine: eight CPU cores,
$30$\,GB of memory, Linux under WSL2, and a consumer GPU that the training did
not use. Reference runtimes measured during the study: one training seed of
one million steps takes about $25$ to $40$ minutes with four parallel
environment processes, so the main run across five seeds takes roughly two and a half
hours and the ablation arm about three hours, including the validation scoring
of every checkpoint. The full evaluation of one configuration over the thirteen
windows takes a few minutes; the boundary sweep's two hundred and sixty
reruns take under an hour.

\subsection{Verification}

Three checks keep the reproductions honest. The accounting identity of
Appendix~\ref{app:system} is asserted on every replay. The parity harness
bounds the distance between the two engines on the recorded tape. The test
suite contains more than two thousand three hundred tests covering the data
boundary, position mathematics, the exact fee engine, the environment,
training, evaluation and the command-line interfaces, at about $95\,\%$ line
coverage. The originality screen compares the report's prose with the cited
papers and the repository's own documents at a threshold of eight consecutive
words; it reports no external matches.
